\subsection{The integral closure of a Krull domain in a finite extension of its field of fractions}

PROPOSITION 12. Let A be a Krull domain, K its field of fractions, K' a finite extension of K and A' the integral closure of A in K'. Then A' is a Krull domain. The essential valuations of A' are the normed discrete valuations on K' which are equivalent to the extensions of the essential valuations of A.

Let \((v_{\iota})_{\iota\in I}\) be the family of extensions to \(K'\) of the essential valuations of A. Since the degree \(n = [K':K]\) is finite, the \(v_{\iota}\) are discrete valuations on \(K'\) (Chapter VI, § 8, no. 1, Corollary 3 to Proposition 1). Let \(B_{\iota}\) be the ring of \(v_{\iota}\) ; then \(A' \subset \bigcap_{\iota\in I}B_{\iota}\) (Chapter VI, § 1, no. 3, Theorem 3). Conversely, every element \(x\) of \(\bigcap_{\iota\in I}B_{\iota}\) is integral over each of the essential valuation rings of A (Chapter VI, § 1, no. 3, Corollary 3 to Theorem 3); hence the coefficients of the minimal polynomial of \(x\) over K belong to A (Chapter V, § 1, no. 3, Corollary to Proposition 11), so that \(x\in A'\) ; thus \(A' = \bigcap_{\iota\in I}B_{\iota}\) . Now let \(x\) be a non-zero element of \(A'\) ; it satisfies an equation of the form \(x^{s} + a_{s-1}x^{s-1} + \cdots + a_{0} = 0\) where \(a_{i}\in A\) and \(a_{0}\neq 0\) ; if \(v_{\iota}(x) > 0\) , then \(v_{\iota}(a_{0}) > 0\) ; now the essential valuations \(v\) of A such that \(v(a_{0}) > 0\) are finite in number and the valuations on \(K'\) extending a given valuation on K are also finite in number (Chapter VI, § 8, no. 3, Theorem 1); hence \(v_{\iota}(x) = 0\) except for a finite number of indices \(\iota\in I\) . Thus it has been proved that \(A'\) is a Krull domain (no. 3, Definition 3).

It remains to show that the \(v_{\iota}\) are equivalent to essential valuations of \(A'\) (no. 4, Corollary 2 to Proposition 6), that is (no. 6, Corollary 2 to Theorem 3) that the prime ideal \(\mathfrak{p}_{\iota}\) , consisting of the \(x \in A'\) such that \(v_{\iota}(x) > 0\) , is of height 1. If this were not so, there would exist a prime ideal \(\mathfrak{q}\) of \(A'\) such that \((0) \subset \mathfrak{q} \subset \mathfrak{p}_{\iota}\) distinct from (0) and \(\mathfrak{p}_{\iota}\) ; then \((0) \subset \mathfrak{q} \cap A \subset \mathfrak{p}_{\iota} \cap A\) and \(\mathfrak{q} \cap A\) would be distinct from (0) and \(\mathfrak{p}_{\iota} \cap A\) (Chapter V, § 2, no. 1, Corollary 1 to Proposition 1); the prime ideal \(\mathfrak{p}_{\iota} \cap A\) would therefore not be of height 1, which contradicts the fact that it corresponds to an essential valuation of \(A\) .

COROLLARY. Let \(\mathfrak{p}\) (resp. \(\mathfrak{p}'\) ) be a prime ideal of A (resp. A') of height 1 and \(v\) (resp. \(v'\) ) the essential valuation of A (resp. A') corresponding to it. For \(\mathfrak{p}'\) to lie above \(\mathfrak{p}\) , it is necessary and sufficient that the restriction of \(v'\) to K be equivalent to \(v\) .

The valuation \(v'\) is equivalent to the extension of an essential valuation w of A (Proposition 12). Let \(q = p' \cap A\) , which is a prime ideal of A of height 1. For the restriction of \(v'\) to K to be equivalent to v, it is necessary and sufficient that w = v and hence that q = p.

